\section{Một phòng thủ, và những con số của nó}
\label{sec:ch14-mot-phong-thu}

\index{SHLIME|(}%
Mục trước đã lập luận rằng thiết lập của \tn{phép dựng giàn}{scaffolding} cho
một điều kiện cần ở mức phương pháp. Bài báo đầu tiên của danh mục mà chương
này đọc lấy đúng thiết lập ấy làm bàn thử nghiệm và hỏi câu tiếp theo: có cách
nào ghép LIME với SHAP để phép tấn công khó thành hơn
không~\autocite{p27shlime}.

\index{SHLIME!bài báo là một bài lặp lại}%
Loại bài phải nói trước, vì nó quyết định cách đọc mọi con số bên dưới. Bài dài
bảy trang, sáu tác giả, và nửa đầu của nó là một phép lặp lại: các tác giả chạy
lại thí nghiệm COMPAS của Slack và cộng sự bằng chính mã nguồn mà các tác giả
kia công bố, và thu được kết quả gần trùng~\autocite{p27shlime}. Bài tự đặt cho
mình phạm vi của một dự án có hạn định thời gian. Nó không có bảng kết quả nào;
mọi phát biểu định lượng trong bài đều là một ngưỡng đọc ra từ một biểu đồ tán
xạ. Đó không phải một lời chê, nhưng nó là điều kiện để đọc phần sau.

\index{SHLIME!hai lỗ hổng bù nhau}%
Quan sát mà bài xuất phát là quan sát về hình dạng của hai đường cong trong
hình của Slack và cộng sự. Trục hoành là điểm F1 của
\tn{bộ phát hiện ngoài phân phối}{out-of-distribution detector},
tức là phép tấn công định tuyến đúng tới đâu; trục tung là tỉ lệ lời giải thích
vẫn còn xếp thuộc tính được bảo vệ ở vị trí thứ nhất, tức là phép tấn công thất
bại tới đâu. Hai đường cong có dáng khác nhau. LIME giữ nguyên mức cao cho tới
một ngưỡng rồi sụp gần như thẳng đứng; SHAP bắt đầu tụt sớm hơn nhiều nhưng tụt
từ từ~\autocite{p27shlime}. Bài đọc hai dáng ấy là hai điểm yếu bù nhau và đề
xuất ghép hai phương pháp lại.

Phép ghép mà bài đưa ra tên là BASIC SHLIME, và nó là một dòng: điểm số của mỗi
đặc trưng là giá trị LIME của đặc trưng ấy nhân với giá trị SHAP của
nó~\autocite{p27shlime}. Lý do bài đưa ra cho phép nhân cũng chỉ một câu, rằng
giá trị SHAP nằm giữa 0 và 1 còn giá trị LIME thì không bị chặn và có thể âm,
nên nhân hai bên lại là cách giữ được nhiều thông tin từ cả hai
nhất~\autocite{p27shlime}.

\index{SHLIME!câu về khoảng giá trị SHAP}%
Câu ấy không đúng, và chương~\ref{ch:shap} là chỗ sách đã nói vì sao. Một
\tn{giá trị Shapley}{Shapley value} có dấu, và độ chính xác cục bộ trong định
nghĩa~\ref{def:ch04-ba-tinh-chat} ràng buộc \emph{tổng} các giá trị bằng hiệu
giữa dự đoán tại $x$ và dự đoán trung bình, chứ không chặn từng giá trị vào
đoạn đơn vị. Chính phần nền của bài, ba trang trước đó, viết mô hình cộng tính
với các hệ số nhận giá trị trên toàn trục số thực~\autocite{p27shlime}, nên bài
tự mâu thuẫn với mình chứ không chỉ nói khác sách. Chỗ này đáng ghi vì nó không
dừng ở một câu viết hớ: cái mà phép nhân làm mất là dấu. Tích của hai số có dấu
mang dấu của tích chứ không mang dấu của thừa số nào, nên một đặc trưng mà LIME
cho điểm âm mạnh và SHAP cũng cho điểm âm mạnh sẽ nhận một điểm số dương lớn,
xếp cùng chỗ với một đặc trưng mà cả hai đều cho điểm dương mạnh. Hai tình
huống ấy là hai tình huống ngược nhau, và điểm số của BASIC SHLIME không phân
biệt được chúng.

\index{SHLIME!ba giá trị cho một ngưỡng}%
Còn các con số thì phải xếp cạnh nhau mới thấy vấn đề, nên xếp chúng ra đây.
Bài phát biểu ngưỡng sụp của LIME ba lần với ba giá trị. Trong đoạn nêu động
cơ, một câu duy nhất viết rằng LIME trụ được cho tới những điểm F1 cao hơn, F1
xấp xỉ 0.7, rồi bị lừa dễ dàng khi vượt ngưỡng ấy, F1 lớn hơn
0.8~\autocite{p27shlime}. Trong đoạn thuật lại phép lặp lại thí nghiệm, cùng
hiện tượng ấy được ghi là tăng vọt khi F1 vượt 0.75~\autocite{p27shlime}.
Ngưỡng khởi phát của SHAP cũng có hai giá trị, F1 xấp xỉ 0.45 ở chỗ này và
khoảng 0.5 ở chỗ kia~\autocite{p27shlime}.

Rồi bài kết luận rằng BASIC SHLIME giữ được tỉ lệ lời giải thích còn xếp đúng
thuộc tính được bảo vệ ở mức trên 0.8 cho một dải F1 rộng hơn và chỉ bắt đầu
suy giảm rõ quanh F1 xấp xỉ 0.75~\autocite{p27shlime}. Con số 0.75 ấy nằm gọn
bên trong khoảng 0.7 tới 0.8 mà chính bài đã dùng cho hiện tượng nó đang so
sánh với. Bài không báo cáo một khoảng tin cậy nào, không báo cáo số lần chạy,
và không có bảng nào để đối chiếu, nên không có cách nào tách cải thiện được
tuyên bố ra khỏi độ tản mạn của chính phép đọc biểu đồ. Sách vì thế ghi nhận
phép dựng và không nhận kết luận về cải thiện.

\index{SHLIME!bài tự nói phép so không công bằng}%
Phần công việc tương lai của bài nói cùng điều ấy bằng lời của chính các tác
giả, và nói mạnh hơn. Câu thứ nhất ghi rằng bước tự nhiên tiếp theo là dựng một
bộ phát hiện ngoài phân phối riêng cho SHLIME, và rằng sức mạnh của các kết quả
hiện có bị hạn chế bởi chỗ không so sánh được độ bền của nó với LIME và SHAP
riêng lẻ trên một thang công bằng~\autocite{p27shlime}. Đó là bài tự phát biểu
rằng phép so mà kết luận của nó đứng trên không phải một phép so công bằng. Còn
phần tóm tắt của bài thì viết rằng nhóm đánh giá nhiều cấu hình ghép và tìm ra
những cấu hình cải thiện đáng kể việc phát hiện thiên lệch, cả hai đều ở số
nhiều~\autocite{p27shlime}; phần thân bài đánh giá đúng một cấu hình, là BASIC
SHLIME. Hai cách ghép khác có được nêu tên, nhưng một cách, lấy trung bình có
trọng số, bị bài loại ngay tại chỗ vì làm mất thông tin của cả hai phía, còn
cách kia thì để lại cho nghiên cứu sau~\autocite{p27shlime}. Không cách nào
được đem ra đo.

\index{SHLIME!không đo cái mà nó phải giữ}%
Câu thứ hai trong phần công việc tương lai mới là câu đáng mang sang
chương~\ref{ch:khoang-trong-nghien-cuu}, và nó khép mục này lại. Các tác giả
ghi rằng nghiên cứu của họ tập trung vào độ bền của SHLIME trước loại tấn công
này, rằng cũng cần xét xem nó làm tốt tới đâu ở chính nhiệm vụ diễn giải mà nó
sinh ra để làm, và rằng việc so SHLIME với LIME và SHAP trên nhiều bộ dữ liệu
là việc của nghiên cứu sau~\autocite{p27shlime}. Nói thẳng ra thì đây là một
phương pháp giải thích được đề xuất, được đo bằng đúng một đại lượng là mức độ
nó chống được một phép tấn công, và chưa từng được đo xem nó có giải thích được
gì hay không.

Hình dạng ấy đã có tên trong sách. Mục~\ref{sec:ch11-dieu-chua-giai-quyet} khép
lại trên một vòng lặp không có bước loại: một phương pháp được đề xuất, được so
với những phương pháp đã có trên một chỉ số do chính người đề xuất chọn, rồi
được công bố. Ở đây vòng lặp ấy quay thêm một lần nữa, và lần này ở nhánh sửa
chữa chứ không ở nhánh đề xuất phương pháp mới. Chỉ số do người đề xuất chọn là
tỉ lệ chống tấn công, và cái bị bỏ ra ngoài phép đo là chính thứ mà công cụ tồn
tại để làm.

\index{SHLIME|)}%
